% Failure-mode / practice coverage matrix as a heat-strip.
% Companion to Table~\ref{tab:failure-modes} in
% paper/sections/failure-modes.tex: where the table reads top-to-bottom,
% this figure reads left-to-right at a glance and exposes which of
% the eight integrated practices bears the load for each named
% failure mode. A solid blue cell marks the practice that primarily
% addresses that mode; a soft cell marks a secondary contributor; an
% empty cell records that the row's failure mode is residual or out
% of scope for that practice. The point of the figure is the shape:
% no row is empty; no practice is unused; the residual row at the
% bottom is intentionally white across all eight columns.
% Loaded from paper/sections/failure-modes.tex via \input.
\begin{figure}[!ht]
\centering
\begin{tikzpicture}[
  every node/.style={font=\sffamily\footnotesize, align=center,
                     inner sep=2pt},
  rowlbl/.style={font=\sffamily\footnotesize, text=black,
                 anchor=east, align=right, text width=42mm},
  collbl/.style={font=\sffamily\bfseries\footnotesize, text=dlrGrau1,
                 align=center, rotate=45, anchor=west,
                 text width=22mm},
  cellprimary/.style={fill=dlrBlau1, draw=dlrGrau4, line width=0.3pt,
                      minimum width=8mm, minimum height=6mm},
  cellsecondary/.style={fill=dlrBlau5, draw=dlrGrau4, line width=0.3pt,
                        minimum width=8mm, minimum height=6mm},
  cellresidual/.style={fill=dlrGelb5, draw=dlrGrau4, line width=0.3pt,
                       minimum width=8mm, minimum height=6mm},
  cellempty/.style={fill=white, draw=dlrGrau4, line width=0.3pt,
                    minimum width=8mm, minimum height=6mm},
  legend/.style={font=\sffamily\scriptsize, text=dlrGrau1},
]
% Practice indices 1..8 across the top, x = i*0.95
\foreach \i/\name in {%
  1/{1.\,Transcripts},
  2/{2.\,Verification\\status (FSM)},
  3/{3.\,Per-claim\\provenance},
  4/{4.\,Mirror\\discipline},
  5/{5.\,Recursive\\meta-process},
  6/{6.\,Base-rate\\disclosure},
  7/{7.\,Legal\\honesty},
  8/{8.\,FAIR as\\precondition}}
{
  \node[collbl] at ({\i*0.95}, 0.35) {\name};
}

% Row layout: y = -row*0.7
% Define rows: label, then which columns are primary (P), secondary (S),
% residual/yellow (R). Order of columns as listed.
% 1 Fabricated citations:           P=2,3   S=6
% 2 Claim-evidence drift:           P=2     S=3,4
% 3 Sloppification (quote drift):   P=4     S=2
% 4 Model collapse (residual):      R=1,6
% 5 List-of-lists prose:            P=5     (readability scrutinizer = practice 5? actually the readability scrutinizer is part of the pipeline; closest practice mapping is 5 recursive meta-process. Use S=5.)  -> use S=5; primary = none -> mark residual on column 5.
% 6 Hedging chains:                 S=5
% 7 Placeholder dishonesty:         P=4
% 8 Dual use:                       R=7   S=3
% Residual stripe at the very bottom: failure modes the pattern only
% signals toward (audience fatigue, adversarial review, leakage).

% Helper: a pair of \node lines per cell.
% Row 1: Fabricated citations
\node[rowlbl] at (0.2, -0.7) {Fabricated citations};
\foreach \c in {1,4,5,7,8} {\node[cellempty]    at ({\c*0.95}, -0.7) {};}
\foreach \c in {2,3}       {\node[cellprimary]  at ({\c*0.95}, -0.7) {};}
\foreach \c in {6}         {\node[cellsecondary]at ({\c*0.95}, -0.7) {};}

% Row 2: Claim-evidence drift
\node[rowlbl] at (0.2, -1.4) {Claim--evidence drift};
\foreach \c in {1,5,6,7,8} {\node[cellempty]    at ({\c*0.95}, -1.4) {};}
\foreach \c in {2}         {\node[cellprimary]  at ({\c*0.95}, -1.4) {};}
\foreach \c in {3,4}       {\node[cellsecondary]at ({\c*0.95}, -1.4) {};}

% Row 3: Sloppification (quotation drift)
\node[rowlbl] at (0.2, -2.1) {Sloppification (quote drift)};
\foreach \c in {1,3,5,6,7,8} {\node[cellempty]    at ({\c*0.95}, -2.1) {};}
\foreach \c in {4}           {\node[cellprimary]  at ({\c*0.95}, -2.1) {};}
\foreach \c in {2}           {\node[cellsecondary]at ({\c*0.95}, -2.1) {};}

% Row 4: Model collapse (indirectly mitigated)
\node[rowlbl] at (0.2, -2.8) {Model collapse (indirect)};
\foreach \c in {2,3,4,5,7,8} {\node[cellempty]    at ({\c*0.95}, -2.8) {};}
\foreach \c in {1,6}         {\node[cellresidual] at ({\c*0.95}, -2.8) {};}

% Row 5: List-of-lists prose
\node[rowlbl] at (0.2, -3.5) {List-of-lists prose};
\foreach \c in {1,2,3,4,6,7,8} {\node[cellempty]    at ({\c*0.95}, -3.5) {};}
\foreach \c in {5}             {\node[cellsecondary]at ({\c*0.95}, -3.5) {};}

% Row 6: Hedging chains
\node[rowlbl] at (0.2, -4.2) {Hedging chains};
\foreach \c in {1,2,3,4,6,7,8} {\node[cellempty]    at ({\c*0.95}, -4.2) {};}
\foreach \c in {5}             {\node[cellsecondary]at ({\c*0.95}, -4.2) {};}

% Row 7: Placeholder dishonesty
\node[rowlbl] at (0.2, -4.9) {Placeholder dishonesty};
\foreach \c in {1,2,3,5,6,7,8} {\node[cellempty]    at ({\c*0.95}, -4.9) {};}
\foreach \c in {4}             {\node[cellprimary]  at ({\c*0.95}, -4.9) {};}

% Row 8: Dual use (made visible, not eliminated)
\node[rowlbl] at (0.2, -5.6) {Dual use (made visible)};
\foreach \c in {1,2,4,5,6,8}   {\node[cellempty]    at ({\c*0.95}, -5.6) {};}
\foreach \c in {7}             {\node[cellresidual] at ({\c*0.95}, -5.6) {};}
\foreach \c in {3}             {\node[cellsecondary]at ({\c*0.95}, -5.6) {};}

% Row 9: Residual band -- pattern can only signal
\node[rowlbl, text=dlrGrau1] at (0.2, -6.5)
  {Residual: leakage, audience fatigue,\\adversarial review};
\foreach \c in {1,...,8}       {\node[cellempty]    at ({\c*0.95}, -6.5) {};}
\draw[draw=dlrGrau3, line width=0.3pt, dashed]
  ({0.95-0.475}, -6.18) -- ({8*0.95+0.475}, -6.18);

% --- Legend ---------------------------------------------------------
\node[legend, anchor=north west, draw=dlrGrau4, fill=white,
      inner sep=4pt, align=left]
  at (-3.6, -7.1)
  {%
    \tikz\draw[fill=dlrBlau1, draw=dlrGrau4]  (0,0) rectangle (3mm,3mm);\,Primary load \quad
    \tikz\draw[fill=dlrBlau5, draw=dlrGrau4]  (0,0) rectangle (3mm,3mm);\,Secondary contribution \quad
    \tikz\draw[fill=dlrGelb5, draw=dlrGrau4]  (0,0) rectangle (3mm,3mm);\,Indirect / residual \quad
    \tikz\draw[fill=white, draw=dlrGrau4]     (0,0) rectangle (3mm,3mm);\,Out of scope for that practice
  };
\end{tikzpicture}
\caption{Coverage of the eight named failure modes by the eight
integrated practices. A primary cell (solid blue) marks the practice
that bears the failure-mode load; a secondary cell (soft blue)
marks a contributing practice; an indirect/residual cell (yellow)
records that the mode is mitigated only at the labelling layer. The
bottom row records the residual failure modes the pattern can only
signal toward (\S\ref{sec:failure-modes}, \emph{Residual}). The
shape of the figure is the argument: no failure-mode row is empty,
no practice column is unused, and the integration --- not any single
practice --- is what produces full coverage.}
\label{fig:failure-mode-coverage}
\end{figure}
